\section{SysADL Co-Simulation and Animation Architecture}

This section describes the co-simulation and visual animation architecture developed for SysADL models. The architecture provides a dual execution pipeline: one optimized for automated, scenario-driven test verification (the simulator), and another designed for real-time visual inspection of the architectural communication (the animator). Both pipelines share a unified semantic foundation—the SysADL Framework—and are powered by a Parsing Expression Grammar (PEG) parser and a dedicated Code Transformer.

---

\subsection{Part I: Conceptual Structure and Solution Architecture}

The simulation framework is organized as a three-stage processing pipeline: parsing, code transformation, and execution. This design separates compiler concerns (syntax checking and target code generation) from execution concerns (runtime environments, scheduling, and visualization).

\subsubsection{Frontend Compilation Pipeline}
The compilation pipeline begins with the \textbf{PEG Grammar} and its generated \textbf{Parser}. The grammar formalizes the rules of the SysADL language, supporting classical structural elements (components, ports, and connectors) and scenario-driven elements (environments, configurations, scenes, timed event injections, and loop structures). The parser reads the raw textual representation of the model and produces a structured Abstract Syntax Tree (AST) represented as a JSON node tree. If the input file violates the grammar rules, the parser immediately raises syntactic exceptions with precise location data (line and column coordinates).

The \textbf{Transformer} acts as the transpilation engine. It processes the AST generated by the parser and generates target code in JavaScript. The transformer maps structural boundaries, connections, and allocations into executable Javascript modules. It isolates compiler logic from the execution environment, producing clean files that represent the system behavior. For advanced environment models, the transformer automatically outputs two files: one for the structural core of the system, and another that specifies the testing environment, activities, and scenarios.

\subsubsection{The Execution Engine Dichotomy}
The execution backend is split into two engines designed for different use cases:
\begin{enumerate}
    \item \textbf{The Simulator (CLI Engine)}: Implemented in \texttt{SysADLSimulator.js}, this component is optimized for automated headless verification. It executes scenario-driven test plans, manages parallel and sequential execution flows, validates pre- and post-conditions of system scenes, and persists comprehensive test reports locally. It is designed to run in continuous integration (CI) pipelines or local CLI shells.
    \item \textbf{The Animator (Web Engine)}: Implemented in \texttt{simulator.js}, this component is designed for interactive, visual inspection. Running entirely client-side in the browser, it loads classical structural models and stimulates boundary ports either through user-defined parameters or loop intervals. It is coupled with the Monaco Editor and a 2D canvas visualization pane to render data transfer packets traveling along connectors, making it an excellent debugging dashboard.
\end{enumerate}

By separating execution into a fast, headless CLI Simulator and a highly visual Web Animator, the architecture provides both rigorous testing capabilities and developer-friendly debugging interfaces using a single, unified parser and transformer code base.

---

\subsection{The SysADL Framework: The Semantic Runtime Foundation}

The core of the simulation architecture is the \textbf{SysADL Framework}, located in the \texttt{sysadl-framework/} directory. The framework defines the execution semantics of the language. It provides the concrete implementations of components, behaviors, and environment definitions that underpins the transpiled code.

\subsubsection{Shared Semantic Foundation}
The framework serves as the unified base that binds compilation, simulator execution, and visual animation together:
\begin{enumerate}
    \item \textbf{Transformer Integration}: The transformer does not output low-level execution routines. Instead, it generates Javascript class definitions that extend the base classes of the framework (such as \texttt{Component}, \texttt{Port}, \texttt{Connector}, and \texttt{EnvComponent}).
    \item \textbf{Simulator Integration}: The CLI simulator loads the transpiled classes and runs execution scenarios. It interacts directly with the framework instances, query state variables via the framework's reactive managers, and propagates event signals through framework-managed connections.
    \item \textbf{Animator Integration}: The web animator runs the compiled code client-side, using the browser-compatible version of the framework. It connects a custom listener (\texttt{SimulationLogger}) to the framework's connector classes to capture signal packet transfers in real time.
\end{enumerate}

\subsubsection{Core Framework Classes}
The framework is divided into structural base classes and scenario/environment base classes defined in \texttt{SysADLBase.js}:
\begin{itemize}
    \item \textbf{\texttt{Component}}: Represents a structural block in the system. It manages its own set of ports and state variables, exposes an evaluation method for its behavioral actions, and triggers delegated behaviors.
    \item \textbf{\texttt{Port}}, \textbf{\texttt{SimplePort}}, and \textbf{\texttt{CompositePort}}: Define boundary connection points. Ports are typed and directional (\texttt{in}, \texttt{out}, or \texttt{inout}). Simple ports transmit atomic value types, while composite ports group multiple ports to represent complex interfaces.
    \item \textbf{\texttt{Connector}}: Manages value propagation between ports. When a port transmits a value, the connector handles the transmission, checks compatibility constraints, and delivers the payload to the receiving port.
    \item \textbf{\texttt{EnvComponent}}, \textbf{\texttt{EnvPort}}, and \textbf{\texttt{EnvConnector}}: Represent mock entities, sensors, actuators, and environmental boundaries that wrap or interact with the system components.
    \item \textbf{\texttt{Scene}}, \texttt{Scenario}, and \texttt{ScenarioExecution}: Define the operational test hierarchy. A Scene checks preconditions and postconditions, waiting for specific event signals. Scenarios group scenes into validation sequences, and ScenarioExecution drives the execution modes.
\end{itemize}

---

\subsection{Part II: Technical Implementation and Mechanics}

This section details the low-level implementation details of the parser, transformer, simulator, and animator components.

\subsubsection{PEG Parser and Grammar Compilation}
The grammar is specified in \texttt{sysadl.peg} using the Parsing Expression Grammar (PEG) formalism. It is compiled into a JavaScript module (\texttt{sysadl-parser.js}) via the Peggy compiler. Unlike Context-Free Grammars, PEG parses text using prioritized choice, ensuring deterministic parsing with linear time complexity ($O(N)$) and eliminating lookahead ambiguity. 

The parser yields an AST composed of structured JSON nodes. For example, a component definition node captures its ports and properties, while scenario executions are parsed with detailed metadata including modes (\texttt{once}, \texttt{loop: N}, or \texttt{loop: while <condition>}) and sequenced statements (such as parallel or repeated execution blocks).

\subsubsection{The Transpilation Engine (transformer.js)}
The transpiler, implemented in \texttt{transformer.js}, compiles the parsed AST JSON structure into standard JavaScript class declarations. A primary design requirement of this transformer is the \emph{preservation of representativeness}—ensuring that every architectural concept defined in SysADL maps directly and cleanly to an equivalent object-oriented structure in JavaScript, preserving naming scopes, hierarchy, and structural topology without semantic loss.

\paragraph{Semantic Mapping and Isomorphism}
To achieve representativeness preservation, the transformer uses prefix mapping to translate SysADL elements into JavaScript class declarations, avoiding collisions and maintaining traceability back to the original model:
\begin{itemize}
    \item \texttt{Component def} $\rightarrow$ Transpiled to classes prefixed with \texttt{CP\_Elements\_} (e.g., \texttt{CP\_Elements\_TempMonitorCP}) inheriting from the framework's \texttt{Component} class.
    \item \texttt{Port def} $\rightarrow$ Transpiled to classes prefixed with \texttt{PT\_Elements\_} (e.g., \texttt{PT\_Elements\_TempIPT}) extending \texttt{SimplePort} or \texttt{CompositePort}.
    \item \texttt{Connector def} $\rightarrow$ Transpiled to classes prefixed with \texttt{CN\_Elements\_} extending \texttt{Connector}.
    \item \texttt{value type} and \texttt{datatype} $\rightarrow$ Mapped to data type descriptors prefixed with \texttt{DT\_} and value wrappers.
\end{itemize}

\paragraph{Code Generation Outputs}
When environment elements are declared in the model, the transformer splits the output into two separate modules:
\begin{itemize}
    \item \texttt{<ModelName>.js}: Exposes the core components, ports, and internal configurations, exporting a factory function named \texttt{createModel()}.
    \item \texttt{<ModelName>-env-scen.js}: Contains the mock environment definitions, scenario scripts, and timing hooks. It exports a factory function named \texttt{createEnvironmentModel()} which imports and wraps the structural model.
\end{itemize}

\paragraph{Examples of Structural Representativeness Preservation}
To illustrate this isomorphism, consider the SysADL port and component definition shown in Listing 1:
\begin{lstlisting}[language=SQL, caption={SysADL Port and Component Definition}]
port def TempIPT { flow in Real }

component def TempMonitorCP {
    ports:
        s1 : TempIPT;
        average : TempOPT;
}
\end{lstlisting}

This compiles to the corresponding JavaScript OOP representation shown in Listing 2, demonstrating how structural relationships are preserved:
\begin{lstlisting}[language=JavaScript, caption={Transpiled JavaScript Port and Component Classes}]
// Transpiled Port Definition
class PT_Elements_TempIPT extends SimplePort {
  constructor(name, opts = {}) {
    super(name, "in", { ...{ expectedType: "Real" }, ...opts });
  }
}

// Transpiled Component Definition
class CP_Elements_TempMonitorCP extends Component {
  constructor(name, opts = {}) {
    super(name, opts);
    this.addPort(new PT_Elements_TempIPT(portName_s1, { owner: name, originalName: "s1" }));
    this.addPort(new PT_Elements_TempOPT(portName_average, { owner: name, originalName: "average" }));
  }
}
\end{lstlisting}

Similarly, SysADL connectors and bindings are mapped to structural data flows. Listing 3 displays a connector definition in SysADL:
\begin{lstlisting}[language=SQL, caption={SysADL Connector Definition}]
connector def FarToCelCN {
    participants:
        ~ f : FTempOPT;
        ~ c : CTempIPT;
    flows: Real from f to c
}
\end{lstlisting}

This is translated to a connector description in JavaScript (Listing 4), maintaining the exact participant and flow semantics:
\begin{lstlisting}[language=JavaScript, caption={Transpiled JavaScript Connector Class}]
class CN_Elements_FarToCelCN extends Connector {
  constructor(name, opts = {}) {
    super(name, {
      participants: {
        f: "PT_Elements_FTempOPT",
        c: "PT_Elements_CTempIPT"
      },
      flows: [
        { type: "Real", from: "f", to: "c" }
      ]
    });
  }
}
\end{lstlisting}

Finally, Listing 5 demonstrates a SysADL ScenarioExecution definition containing mode configuration and dynamic signal injections:
\begin{lstlisting}[language=SQL, caption={SysADL ScenarioExecution Definition}]
ScenarioExecution RobAFIS_Execution to ValidationPlan {
    mode: once;
    inject StartSimulationSig { mission = P0; } immediate;
    Scenario_RoutingLogic;
}
\end{lstlisting}

During transpilation, the transformer compiles this execution statement into a sequenced asynchronous subroutine (Listing 6):
\begin{lstlisting}[language=JavaScript, caption={Transpiled JavaScript ScenarioExecution Class}]
class SE_RobAFIS_Execution extends ScenarioExecution {
  constructor() {
    super("RobAFIS_Execution", "ValidationPlan", "once");
  }
  
  async executeAsync(ctx) {
    // inject StartSimulationSig { mission = P0; } immediate;
    const injectData = { StartSimulationSig: { mission: "P0" } };
    if (ctx.envActivities) {
      ctx.envActivities.handleSignal('StartSimulationSig', injectData.StartSimulationSig, ctx);
    }
    // Scenario_RoutingLogic;
    await this.executeScenario('Scenario_RoutingLogic', ctx);
  }
}
\end{lstlisting}

As shown, the transformer maintains structural integrity, typing constraints, naming hierarchies, and behavioral sequences. No architectural meaning is lost in translation, facilitating seamless semantic co-simulation in both CLI and browser animators.

\subsubsection{SysADLSimulator.js: The Scenario CLI Simulator}
The simulator is a unified CLI executor. Its execution flow consists of the following steps:
\begin{enumerate}
    \item \textbf{Context Proxying}: It creates a dynamic context proxy using a Javascript \texttt{Proxy} pattern. This proxy intercepts read and write operations to model properties, allowing the simulator to track changes in component state variables and route delegations through boundary extensions.
    \item \textbf{Simulation Scheduler}: A timing engine called \texttt{SimulationScheduler} resolves chronological delays (\texttt{after 5s}), scenario boundaries (\texttt{before/after scenario}), and reactive triggers (\texttt{when <condition>}). It uses a polling scheduler combined with state change watchers to evaluate conditional guards.
    \item \textbf{Replication Tracking}: To handle replicated component collections (such as arrays of parts), the proxy walks up the component parent pointers (\texttt{.parent}) to locate the closest replicated ancestor, incrementing indices in a commit-phase manner to prevent index corruption.
    \item \textbf{JUnit Reporting and JSONL Logging}: After completing scenario trials, the simulator evaluates the pass/fail status of pre- and post-conditions, formats a JUnit-style log report to the terminal, and writes execution data into JSON Lines (\texttt{.jsonl}) files under the \texttt{./logs} directory.
\end{enumerate}

\subsubsection{simulator.js: The Web Animator}
The web animator is designed to execute code in the browser sandbox. Because browsers restrict raw file system access (CORS), the animator runs entirely in memory:
\begin{itemize}
    \item \textbf{Dynamic Code Evaluation}: The web application orchestrator (\texttt{app.js}) fetches the transpiled Javascript code from a local Node server endpoint (\texttt{/api/transform}) and evaluates the string using an IIFE wrapper:
    \begin{lstlisting}[language=JavaScript]
    (function() {
      var module = { exports: {} };
      var exports = module.exports;
      // Transpiled code
      return module.exports;
    })()
    \end{lstlisting}
    \item \textbf{Client-side Module Resolution}: The animator mocks the Node.js \texttt{require()} system inside the browser, intercepting framework paths and mapping them to browser globals:
    \begin{lstlisting}[language=JavaScript]
    function require(p) {
      if (p.includes('SysADLBase')) return window.SysADLBase;
      if (p.includes('SimulationLogger')) return window.SimulationLogger;
      // Return stubs for Node-specific modules (e.g. TaskExecutor)
      return { TaskExecutor: class {} };
    }
    \end{lstlisting}
    \item \textbf{Visual Trace Integration}: When components communicate, the framework calls methods in \texttt{SimulationLogger.js}. The logger records these actions to a list of events:
    \begin{lstlisting}[language=JSON]
    {
      "timestamp": 64,
      "flowId": "flow_1",
      "type": "CONNECTOR_TRIGGERED",
      "data": { "connectorName": "c1", "from": "s1.temp", "to": "tempMon.s1" }
    }
    \end{lstlisting}
    A canvas-based graphics engine (\texttt{visualizer.js}) reads these logged events, computes the coordinate points of component nodes and connector paths, and plays smooth animations of data packets flowing between inputs and outputs.
\end{itemize}

\subsubsection{Reactive Monitoring Integration}
To handle reactive monitoring for classical models in the browser, the animator integrates with the framework's reactive watcher. Since \texttt{ReactiveConditionWatcher} does not extend an event emitter class, the simulator dynamically inspects the watcher's method structure:
\begin{lstlisting}[language=JavaScript]
if (typeof watcher.emit === 'function') {
  // Bind to event-based emitter
} else if (typeof watcher.triggerCondition === 'function') {
  // Intercept method call directly (Reactive System)
  const orig = watcher.triggerCondition.bind(watcher);
  watcher.triggerCondition = function(id, cond, state) {
    console.log(`[REACTIVE] Condition met: ${id}`);
    return orig(id, cond, state);
  };
}
\end{lstlisting}
This wraps the reactive trigger checks, enabling condition logging and visual animations in both Node and browser environments.
